% ============================================================
\chapter{Point Cloud Representations}
\label{ch:point_clouds}
% ============================================================

In 2017, Charles Qi and collaborators at Stanford published PointNet, a neural network architecture capable of processing 3D point clouds directly --- without converting them to voxels or meshes. The key idea is elegant: since a point cloud is an \emph{unordered} set, the architecture must be \emph{permutation invariant}. PointNet achieves this by applying a shared function to each point and then aggregating via a global max. PointNet++, Dynamic Graph CNN (DGCNN), and transformer-based methods subsequently extended this approach by capturing local structures and neighbourhood relationships. This chapter explores these architectures, from theoretical foundations (Zaheer's theorem on set functions) to applications in autonomous driving, robotics, and 3D reconstruction.

\begin{intuition}
A 3D point cloud\index{point cloud} is an \textbf{unordered} set
of points $\{x_1, \ldots, x_n\} \subset \R^3$. Unlike images
(regular grids) or graphs (combinatorial structure), it has no
canonical connectivity. The challenge is to design architectures
that are \textbf{permutation invariant} and, ideally,
\textbf{equivariant} with respect to geometric transformations
(rotations, translations).
\end{intuition}

% ============================================================
\section{The permutation invariance problem}
% ============================================================

\begin{definition}[Permutation-invariant function]
\index{permutation invariance}
A function $f : \R^{n \times d} \to \R^k$ is
\textbf{permutation invariant} if for every permutation matrix
$\Pi \in \{0,1\}^{n \times n}$:
\[
  f(\Pi X) = f(X).
\]
\end{definition}

\begin{theorem}[Characterization --- Zaheer et al.\ 2017]
\label{thm:deepsets}
\index{Deep Sets}
Any continuous permutation-invariant function
$f : 2^{\R^d} \to \R$ can be written as:
\[
  f(\{x_1, \ldots, x_n\}) = \rho\!\left(
  \sum_{i=1}^n \phi(x_i)\right)
\]
for continuous functions $\phi : \R^d \to \R^q$ and
$\rho : \R^q \to \R$ (provided $q$ is sufficiently large).
\end{theorem}

\begin{remark}
This theorem underpins the \textbf{Deep Sets} architecture and,
by extension, \textbf{PointNet}. The key operation is sum (or max)
aggregation, which guarantees permutation invariance.
\end{remark}

% ============================================================
\section{PointNet}
\label{sec:pointnet}
% ============================================================

\begin{definition}[PointNet architecture --- Qi et al.\ 2017]
\index{PointNet}
\textbf{PointNet} directly processes a point cloud of $n$ points
$X \in \R^{n \times 3}$:
\begin{enumerate}
  \item \textbf{Spatial transformer}: a mini-network (T-Net)
        predicts a matrix $T \in \R^{3 \times 3}$ to align the
        cloud: $X' = X \cdot T$.
  \item \textbf{Shared MLP}: each point is transformed
        independently by an MLP $\phi : \R^3 \to \R^{1024}$:
        $h_i = \phi(x_i')$.
  \item \textbf{Symmetric aggregation}: global max-pooling
        produces a global feature vector:
        $g = \max_{i=1}^n h_i$.
  \item \textbf{Classification/segmentation head}: a final MLP
        on $g$ (or $[g; h_i]$ for segmentation).
\end{enumerate}
\end{definition}

\begin{keyformulas}[title=\textbf{PointNet --- fundamental equation}]
\[
  f(X) = \rho\!\left(\max_{i=1}^n \phi(x_i)\right)
\]
where $\phi$ is a pointwise MLP and $\rho$ a global MLP.
\end{keyformulas}

\begin{proposition}[Universal approximation of PointNet]
PointNet can approximate any continuous permutation-invariant
function on compact point clouds, in the sense of
Theorem~\ref{thm:deepsets}.
\end{proposition}

\begin{attention}[title=\textbf{Limitations of PointNet}]
PointNet processes each point \textbf{independently} before
aggregation: it does not capture local structures (neighborhoods).
Two point clouds with the same points but different local
geometries may yield the same global vector.
\end{attention}

% ============================================================
\section{PointNet++}
\label{sec:pointnetpp}
% ============================================================

\begin{definition}[PointNet++ architecture --- Qi et al.\ 2017]
\index{PointNet++}
\textbf{PointNet++} introduces a hierarchy of local processing:
\begin{enumerate}
  \item \textbf{Sampling}: select $n'$ centers via farthest point
        sampling (FPS).
  \item \textbf{Grouping}: for each center $c_j$, select the $K$
        nearest neighbors (or ball query of radius $r$).
  \item \textbf{Local PointNet}: apply PointNet on each group,
        producing a feature $g_j \in \R^{d'}$.
  \item \textbf{Iteration}: repeat on the centers with their
        features, creating a multi-scale pyramid.
\end{enumerate}
\end{definition}

\begin{definition}[Set Abstraction Layer]
\index{Set Abstraction}
The \textbf{Set Abstraction} (SA) layer combines the three
operations:
\[
  \mathrm{SA}(X) = \{g_j\}_{j=1}^{n'}, \qquad
  g_j = \max_{x_i \in \mathcal{N}(c_j)} \phi(x_i - c_j)
\]
where $\mathcal{N}(c_j)$ is the neighborhood of center $c_j$.
\end{definition}

\begin{intuition}
PointNet++ is to point clouds what CNNs are to images: it applies
local operations hierarchically. The difference is the absence of
a grid: neighborhoods are determined dynamically.
\end{intuition}

% ============================================================
\section{Point cloud convolutions}
\label{sec:convolutions}
% ============================================================

\begin{definition}[Continuous convolution on point clouds]
\index{continuous convolution}
A \textbf{continuous convolution} on a point cloud is:
\[
  (f * g)(x_i) = \sum_{x_j \in \mathcal{N}(x_i)}
  g(x_j - x_i) \cdot h_j
\]
where $g : \R^3 \to \R^{d' \times d}$ is a continuous filter
and $h_j \in \R^d$ are the features of point $x_j$.
\end{definition}

\begin{example}[Convolution-based architectures]
\begin{itemize}
  \item \textbf{PointConv} (Wu et al.\ 2019): $g$ is parameterized
        by an MLP applied to relative position and local density.
  \item \textbf{KPConv} (Thomas et al.\ 2019): $g$ is defined
        by fixed kernel points\index{KPConv} $\{y_k\}_{k=1}^K$
        with correlation functions:
        $g(x) = \sum_k W_k \cdot \max(0, 1 - \norm{x - y_k}/\sigma)$.
  \item \textbf{DGCNN} (Wang et al.\ 2019): convolution on a
        dynamic $k$-nearest neighbor graph in feature space.
\end{itemize}
\end{example}

% ============================================================
\section{Rotation equivariance}
\label{sec:rotation_equivariance}
% ============================================================

\begin{definition}[$\mathrm{SO}(3)$ equivariance]
\index{SO(3) equivariance}
A function $f : \R^{n \times 3} \to \R^{n \times d}$ is
\textbf{$\mathrm{SO}(3)$-equivariant} if for every rotation
$R \in \mathrm{SO}(3)$:
\[
  f(R \cdot X) = \rho(R) \cdot f(X)
\]
where $\rho(R)$ is a representation of $R$ acting on the output
space.
\end{definition}

\begin{theorem}[Equivariant convolutions --- Thomas et al.\ 2018]
\index{Tensor Field Networks}
Equivariant filters under $\mathrm{SO}(3)$ are expressed in the
spherical harmonics basis:
\[
  g^{(\ell)}(r, \hat{x}) = R^{(\ell)}(r) \cdot Y^{(\ell)}(\hat{x})
\]
where $Y^{(\ell)}$ are spherical harmonics of degree $\ell$,
$R^{(\ell)}$ is a learned radial function, and
$\hat{x} = x / \norm{x}$.
\end{theorem}

\begin{remark}
Equivariant architectures (Tensor Field Networks, SE(3)-Transformers,
EGNN) guarantee that predictions transform correctly under rotation.
This eliminates the need for data augmentation with random rotations.
\end{remark}

% ============================================================
\section{Applications}
\label{sec:applications_pc}
% ============================================================

\begin{example}[3D object classification]
\index{3D classification}
On the \textbf{ModelNet40} benchmark (12,311 models, 40 classes):
\begin{center}
\begin{tabular}{lc}
\toprule
\textbf{Method} & \textbf{Accuracy (\%)} \\
\midrule
PointNet & 89.2 \\
PointNet++ & 91.9 \\
DGCNN & 92.9 \\
KPConv & 92.9 \\
Point Transformer & 93.7 \\
\bottomrule
\end{tabular}
\end{center}
\end{example}

\begin{example}[Semantic scene segmentation]
\index{semantic segmentation}
For LiDAR point cloud segmentation (S3DIS, ScanNet), PointNet++
and KPConv achieve competitive performance. Segmentation assigns
each point a semantic label (floor, wall, furniture, etc.).
\end{example}

\begin{minted}{python}
import torch
from torch_geometric.nn import PointNetConv, fps, radius

class PointNetPPLayer(torch.nn.Module):
    """A Set Abstraction layer from PointNet++."""
    def __init__(self, ratio, r, nn):
        super().__init__()
        self.ratio = ratio
        self.r = r
        self.conv = PointNetConv(nn)

    def forward(self, x, pos, batch):
        # Farthest Point Sampling
        idx = fps(pos, batch, ratio=self.ratio)
        # Ball query
        row, col = radius(
            pos, pos[idx], self.r, batch, batch[idx]
        )
        edge_index = torch.stack([col, row], dim=0)
        # Local PointNet
        x_out = self.conv(x, (pos, pos[idx]), edge_index)
        return x_out, pos[idx], batch[idx]
\end{minted}

% ============================================================
\section{Exercises}
% ============================================================

\begin{exercise}
Show that global max-pooling is permutation invariant. Give a
counterexample showing that concatenation is not.
\end{exercise}

\begin{exercise}
Compute the number of parameters in PointNet for a cloud of
$n = 1024$ points in 3D, with a pointwise MLP
$[3, 64, 128, 1024]$ and a classifier $[1024, 512, 256, 40]$.
\end{exercise}

\begin{exercise}[Equivariance]
Show that PointNet with T-Net is \textbf{approximately}
rotation invariant but not exactly equivariant. What property
is missing?
\end{exercise}

\begin{exercise}[Implementation]
Implement a simple PointNet in PyTorch and train it on ModelNet10.
Compare with and without the T-Net.
\end{exercise}

\begin{exercise}[Point cloud convolution]
Explain why KPConv uses fixed kernel points rather than an MLP
to define the filter. What advantage does this provide in terms
of efficiency and interpretability?
\end{exercise}
